\section{トーラス結び目の表現 [1]}\label{sec:torus}

\subsection{入力と h/g 分解}
入力は \cite{AgSc12} §6 の $T(2,2m+1)$ の表現 [1] の一般式（unknot で正規化）である。
$m = 1$ で \cite{AgSc12} (6.15) に一致する。
SO(4) $=$ SU(2)$^2$、SO(6) $=$ SU(4) の同一視で、(6.15) は \cite{GN13} の構成による U$(n)$ の値と一致する（有理点 2 つで厳密）。
$t = q$ では Mariño の Kauffman 多項式の鏡像に、framing の単項式を除いて一致する。

% [T1]
\begin{hypothesis}[h/g 分解]\label{hyp:hg}
$f_1(K) = -W_{\SO}(K) = \tfrac12 h_{\yone} + g_{\yone}$ と書き、被覆の部分を
\[
\tfrac12 h_{\yone} = \mathcal{P}(\ba_{\mathrm{KN}}, \bq, \bt)\,\frac{\lambda - \lambda^{-1}}{\bq - \bq^{-1}},\qquad \ba_{\mathrm{KN}}^2 = a\,(t/q)^{1/2}
\]
と置く。$\mathcal P$ は \cite{DGR05} の superpolynomial で、A 型の refined CS の表現 [1] に当たる。
\end{hypothesis}

% [T2]
\begin{compresult}
$m = 1, \dots, 5$ で、$g_{\yone}$ が Laurent になる単項式 $\pm q^{c/2} t^{d/2} (t^N)^e$ を探索した。条件を満たすのは $e = -2(2m+1)$、$c = -(3m+2)$ に限られる。
$d$ は「$t = q$ の極限が Kauffman 多項式の形」で一意に $d = 7m+4$ と決まる。
したがって単項式は $T_{\yone}^{-2(2m+1)}(q/t)^{\delta/2}$ である（$T_{\yone} = q^{1/2}t^{N-1}$、$\delta = m$）。
一様な規格化では、余分な $(q/t)^{\delta/2}$ は被覆側に移り、SO 側の単項式はちょうど $T_{\yone}^{-2(2m+1)}$ になる（固定 $N = 2, 3$ で data と全係数一致、$m \le 3$）。
\end{compresult}
同じ形の因子 $\sqrt{q/t}$ は、\cite{AgSh11} の SU$(N)$ の trefoil でも現れる。

\subsection{結果}
一様な規格化でのクロスキャップ部分を $g^{\mathrm{true}}_{\yone} = (t/q)^{\delta/2} g_{\yone}$ と書く。
% [T3][T4]
\begin{compresult}\label{cr:torus}
$m = 1, \dots, 5$ で、$g^{\mathrm{true}}_{\yone}$ は整数係数の Laurent 多項式（項数 11, 27, 49, 77, 111）で、次が成り立つ。
\begin{enumerate}
\item[(N)] SO(2) の点 $\ba^2 = \bq^2$（$\lambda^2 = t$）で $g^{\mathrm{true}}_{\yone} = 1$。
\item[(P)] $\bt \to -\bt$ で不変。
\item[($c{=}2$)] $\ba$ の奇数冪の部分は
\[
-(\ba^2 - \bq^2)(\bt^2 - 1)\,\ba^{-(4m+1)}\bq^{-(2m-1)}\bt^{-4m}\,\Pi_m,\qquad
\Pi_m = \sum_{j=0}^{m-1}\ba^{2j}\bq^{2(m-1-j)}\bt^{2j}\sum_{i=0}^{j}(\bq^4\bt^2)^i .
\]
$\Pi_m$ は係数がすべて $+1$ の $m(m+1)/2$ 項から成る。閉じた形は $m \le 3$ から推測し、$m = 4, 5$ で確かめた（予想）。
\item[($c{=}1$)] $\ql$ の整数冪の部分（$\ql^{1/2}$ の偶数冪）は $\sigma$ で対称で、$G$ と $\qr^{1/2}$ の多項式に書ける。$\lambda^\beta$ ごとのブロックは次のどれかである。
  \begin{itemize}
  \item $\lambda^{-(4m+2)}$：正の単項式
  \item $\lambda^{-4m}, \dots, \lambda^{-(2m+4)}$：$(\bt^2 - 1)\times$ 正
  \item $\lambda^{-(2m+2)}$：係数がすべて負
  \item $\lambda^{-2m}$：係数がすべて正
  \end{itemize}
\end{enumerate}
\end{compresult}
$t = q$ では $z$ の奇数冪が消え、トーラス結び目の表現 [1] で $N^{c=2} = 0$ になることと整合する。
$m = 1, 2, 3$ の値は\ref{app:torus} に載せた。

\subsection{判定式}\label{sec:criterion}
unrefined のクロスキャップの BPS 数は整数だが、符号は $\beta$ によって変わる（例：Mariño の Hopf の $\hat g_{\yone,\yone} = z(\lambda - \lambda^{-1})$）。
refined では正値性を期待し、それを「各 $\lambda^\beta$ のブロックの係数が、unrefined の符号 $\varepsilon_\beta$ と同じ符号にそろう」こととして定式化する。

\begin{definition}\label{def:c1}
表現 [1] のクロスキャップ部分の $c = 1$ 部分を
\[
g^{c=1} = \sum_\beta \lambda^\beta \sum_{g, J} N^{c=1}_{g,\beta,J}\, G^g \qr^{2J}
\]
と展開する。$t = q$ で $\sum_J N^{c=1}_{g,\beta,J}$ は Mariño の $N^{c=1}_{g,\beta}$ に戻る。
\end{definition}

% [T6]
\begin{conjecture}[判定式]\label{conj:criterion}
2 ストランドのトーラス結び目 $T(2,2m+1)$ の表現 [1] では、次が成り立つ。
\begin{enumerate}
\item[($c{=}1$)] $G$ 展開の後、$\lambda^\beta$ ごとのブロックは同符号か、$(\bt^2 - 1)\times$ 同符号になり、その符号は unrefined の符号 $\varepsilon_\beta$ である。
\item[($c{=}2$)] $c = 2$ 部分の refined だけの部分 $g^{c=2} - g^{c=2}|_{t=q}(\lambda \to \tlam)$ は $(\ba^2 - \bq^2)(\bt^2 - 1)\times$（$\ba$ のブロックごとに同符号の多項式）である。$T(2,2m+1)$ では unrefined の $c = 2$ 不変量が 0 なので $c = 2$ 部分全体が refined だけの部分で、$-(\ba^2 - \bq^2)(\bt^2 - 1)\times$ 正の多項式になる。
\end{enumerate}
計算結果 \ref{cr:torus} の (N)(P) は規格化の検査で、正値性の内容は ($c{=}1$)($c{=}2$) にある。
\end{conjecture}
根拠は $T(2,2m+1)$（$m \le 5$）の表現 [1] である。結び目一般の予想にはしない：$T(3,4)$ では $c = 1$ 部分が $\sigma$ 対称でなく ($c{=}1$) が成り立たない（($c{=}2$) は成り立つ）。代数的なケーブルでは $c = 1$ 部分が $\sigma$ 対称でなく、$c = 2$ 部分の refined だけの部分は普遍因子で割り切れるが商の符号が $\ba$ のブロック内で混在する（§\ref{sec:cable}）。なお、ケーブルでは unrefined の $c = 2$ 不変量が 0 でないので、$c = 2$ 部分全体が $(\bt^2 - 1)$ で割り切れないことは判定式の破れではない。
$(\ba^2 - \bq^2)$ は SO(2) の点（可換群では結び目の不変量が自明になる）で消える因子である。
$(\bt^2 - 1)$ は「$t = q$ で消える」ことと、(P) の偶奇の帰結である。
